%% CPP 2027 Submission — 0pat: Composite Numbers as Prime Factorization
\documentclass[sigplan,10pt,anonymous,review]{acmart}
\settopmatter{printfolios=true,printacmref=false}

\AtBeginDocument{%
  \providecommand\BibTeX{{\normalfont B\scshape ib\TeX}}}

\settopmatter{printacmref=false}
\renewcommand\footnotetextcopyrightpermission[1]{}
\pagestyle{plain}

\usepackage{microtype}
\usepackage{booktabs}
\usepackage{amsmath,amssymb}

\begin{document}

\title{Composite Numbers as Prime Factorization:\\ A Pure-Known-Mathematics
Re-formalization of a Pat Observation}

\author{Anonymous}
\affiliation{%
  \institution{Anonymous}
}
\email{anonymous@example.org}

\begin{abstract}
We report the first completed exercise of a re-formalization program: a pat
framework observation (``composite numbers are multiple projection
superpositions of primes'') is re-stated in pure known mathematics and
formally proved in Lean 4 / mathlib with \emph{no} pat framework concepts.
The re-statement is the fundamental theorem of arithmetic; the proof is four
theorems: existence of the prime factorization
(\texttt{prod\_factorization\_pow\_eq\_self}), primality of the
factorization support, uniqueness of the factorization (via mathlib's
\texttt{factorizationEquiv}), and the existence of a nontrivial prime
factor for every composite number (\texttt{exists\_prime\_and\_dvd}). All
four compile with zero errors against mathlib v4.32.2. The exercise
demonstrates the program's method: walk the genealogy map from the pat
concept to its standard-mathematics coordinate, formalize the equivalence
on the way, and keep the pat vocabulary out of the proof.
\end{abstract}

\keywords{Lean, formalization, arithmetic, fundamental theorem of arithmetic, 0pat, reproducibility}

\maketitle

\section{Introduction}

\emph{Program.} This paper is the first completed exercise of a
re-formalization program (``0pat''): take the exercises of the pat
framework --- a body of Lean-formalized observations about basepoint-
relative operator semantics --- and re-prove their mathematical content
using \emph{pure known mathematics}: only mathlib's existing definitions
and theorems, no pat framework concepts (basepoints, phase locking,
interlock pairs, projection superpositions).

\emph{Method.} Each exercise is located on a genealogy map (rows: pat
concepts R1--R25; columns: mathematical genealogies --- geometry, algebra,
analytic number theory, topology, category theory, representation,
physics, construction). The map's entries are the paths: a pat concept at
coordinate $(R_i, C_j)$ has a standard-mathematics correspondent at that
coordinate. The re-formalization walks the path: re-state the pat claim at
its coordinate in standard mathematics, formalize the equivalence, and
prove. Since the mathematical claim is known (the ``answer'' exists in
mathlib), the work is path-finding, not theorem-finding.

\emph{This exercise.} The pat observation ``composite numbers are multiple
projection superpositions of primes'' sits at $(R23, C3)$ on the genealogy
map (prime circle $\times$ analytic number theory). Its standard-
mathematics correspondent is the fundamental theorem of arithmetic. We
re-state and prove it.

\section{Re-statement (0pat)}

\paragraph{Pat claim.} A composite number is a superposition of prime
projections: its prime structure is fully determined by the primes that
divide it and their multiplicities.

\paragraph{0pat re-statement (fundamental theorem of arithmetic).}
\begin{enumerate}
\item \emph{Existence:} every nonzero natural number equals the product of
its prime powers: $\prod_{p} p^{e_p} = n$ where $e_p =
\mathrm{factorization}(n)(p)$.
\item \emph{Primality of support:} $e_p > 0$ only for prime $p$.
\item \emph{Uniqueness:} any prime-supported factorization whose product
equals $n$ is exactly $n$'s factorization.
\item \emph{Composite structure:} every composite $n > 1$ has a prime
factor $p < n$.
\end{enumerate}

\section{Formalization}

The proof is a single Lean file (\texttt{proof.lean}, 90 lines), importing
only \texttt{Mathlib.Data.Nat.Factorization.Basic} and
\texttt{Mathlib.Data.Nat.Basic}. Four theorems:

\begin{center}
\begin{tabular}{lll}
\toprule
Theorem & Statement & mathlib anchor \\
\midrule
\texttt{factorization\_represents} & $n \neq 0 \Rightarrow \prod_p p^{e_p} = n$
  & \texttt{Nat.prod\_factorization\_pow\_eq\_self} \\
\texttt{factorization\_support\_prime} & $p \in \mathrm{supp}\, e \Rightarrow$ prime $p$
  & \texttt{Nat.factorization\_eq\_zero\_iff} \\
\texttt{factorization\_unique} & prime support + product $= n \Rightarrow f = n.factorization$
  & \texttt{Nat.factorizationEquiv} \\
\texttt{composite\_has\_prime\_factor} & $1 < n \wedge \neg$prime $n \Rightarrow \exists p$, prime $p$, $p \mid n$, $p < n$
  & \texttt{Nat.exists\_prime\_and\_dvd} \\
\bottomrule
\end{tabular}
\end{center}

\paragraph{Uniqueness via the factorization equivalence.} The cleanest
route to uniqueness is mathlib's
$\mathrm{factorizationEquiv} : \mathbb{N}^+ \simeq \{ f : \mathbb{N}
\rightarrow_0 \mathbb{N} \mid \forall p \in \mathrm{supp}\, f, \text{prime } p \}$:
the inverse maps $f$ to $f.\mathrm{prod}\,(\cdot^\cdot)$ (a mathlib
theorem), so a prime-supported $f$ with product $n$ is the preimage of
$n$ under the inverse --- hence equals $\mathrm{factorizationEquiv}(n)$,
which is definitionally $\langle n.\mathrm{factorization}, \cdot\rangle$.
The whole proof is: $f = \mathrm{factorizationEquiv}(n) = n.\mathrm{factorization}$.

\paragraph{Verification.} Compiled with Lean 4.32.2 (mathlib v4.32.2),
zero errors, zero \texttt{sorry}, zero axioms beyond mathlib's. Build
command in the artifact.

\section{Relation to the pat exercise}

The original exercise (pat\_pnp exercise XXIII, claim R234) formalized the
same observation inside the pat framework, using pat vocabulary. This
re-formalization shows that the mathematical content survives the removal
of the framework: the observation is the fundamental theorem of
arithmetic, and its proof is entirely within known mathematics. The pat
vocabulary is preserved in the exercise record (README) as provenance, not
in the proof.

\section{Conclusion}

The first 0pat exercise is complete: composite numbers as prime
factorization, proved with pure known mathematics in four theorems, zero
errors. The method (genealogy-map path-walking, formalize the equivalence
on the way) is validated on this exercise and is the template for the
remaining ones.

\bibliographystyle{ACM-Reference-Format}
\bibliography{refs}

\end{document}
